% Nouveau EGA V (cote 25): archive cote 25, site page 25 (PDF page 26).
% typescript. Draft compared: none.
% AI-transcribed by Claude Opus 5.5 (claude-opus-5-5, adaptive thinking, effort high), Message Batches API, batch ega25b (msgbatch_01A7n5YuXubUPNK1zsepi8RG), 2026-11-13; not yet checked by eye.
\scanpage{25/25}
\mspage{\hand{21 bis}}

\deleted{ce à peu près immédiate de 15.2.2. appliqué à $\underline{Y} \longrightarrow \check{P}$ .}

\deleted{Il faut voir si  le fait de mélanger dans un même énoncé des proprié-}%
\deleted{tés modulées et non modulables  ne le\xout{s} rend pas  impigeable\xout{s} pratique-}%
\deleted{ment. Noter aussi  que 4.6. et 4.7. est entièrement indépendant des}
\deleted{autres développements des n\textsuperscript{o}s 2,3,4 ; à ce titre, ils méritent sans dou-}%
\deleted{te d'être mis dans un n\textsuperscript{o} à part, où il n'est pas question des propriétés}
\deleted{de la section hyperplane générique, mais uniquement de déterminer si}
\deleted{l'ensemble E est ouvert. On pourrait donner comme titre d\over{a}{e} n\textsuperscript{o} :}
\deleted{``\uline{Théorèmes du type \xout{\illegible{1}} Seidenberg pour une section hyperplane variable}''.}
\deleted{Il faut signaler que le cas  ``géométriquement normal'' (\deleted{\unsure{et}} bien entendu}
\deleted{dans le cas où S=Spec(k) )  est en effet dû à Seidenberg. N'ayant pas}
\deleted{lu son papier, je ne sais s'il part d'une variété X déjà supposée nor-}%
\deleted{male, ou même lisse.  Ce dernier cas semble techniquement nettement}
\deleted{moins vache, puisqu'il suffit de 4.6. qui ne s'appuie pas sur le}
\deleted{délicat 15.2.2.  (Si Seidenberg ne se borne pas à ce cas, je me demande}
\deleted{comment il a pu faire pour se débrouiller).  Si on fait le scindage en}
\deleted{deux envisagé pour le présent numéro, la première partie  devrait prend-}%
\deleted{re le ti\unsure{tre}:  \uline{Etude d'une section hyperplane variable}: \uline{cas des}}
\deleted{\uline{sections hyperplanes} ``\uline{assez générales}''.}

\over{\unsure{6}}{5}. \uline{Connexion d'une section hyperplane quelconque}.

Nous allons ici combiner le critère déjà connu de connexion
géométrique de\xout{xx} la section hyperplane générique  (3.3.), avec le
th de connexion de Zariski, pour obtenir un résultat de connexion pour
une section hyperplane quelconque:
\uline{Proposition} \over{\unsure{5}}{6}.1.  On suppose S=Spec(k), k un corps\xout{xxxx} alg clos, \inserted{X propre sur S} \xout{xx}
que pour toute composante irréductible X$_{i}$ de X, \xout{xxxxxx} $\overline{f(X_{i})}$ soit de
dimension \over{$>$}{$\geqslant$} 2, \xout{xxxxxxxxxxxxxxxxxxxxxxxxxxxxxxxhyper} \deleted{Alors pour}

\margin{25}
